\documentclass[10pt,a4paper]{article}
\usepackage[utf8]{inputenc}
\usepackage[russian,english]{babel}
\usepackage[T2A]{fontenc}
\usepackage{geometry}
\usepackage{enumitem}
\usepackage{hyperref}
\usepackage{titlesec}
\usepackage{xcolor}

% Настройки страницы
\geometry{left=1.8cm, right=1.8cm, top=1.5cm, bottom=1.5cm}
\pagestyle{empty}
\setlength{\parindent}{0pt}

% Цвета
\definecolor{linkcolor}{RGB}{0,102,204}
\hypersetup{
    colorlinks=true,
    linkcolor=linkcolor,
    urlcolor=linkcolor,
}

% Форматирование заголовков
\titleformat{\section}{\large\bfseries}{}{0em}{}[\titlerule]
\titlespacing*{\section}{0pt}{6pt}{3pt}

\begin{document}

% ====== HEADER ======
\begin{center}
    {\Huge \textbf{Хамидуллин Амир}} \\[3pt]
    \small
    \href{mailto:ahamidullin80@gmail.com}{ahamidullin80@gmail.com} | 
    +7 985 138-53-83 | 
    \href{https://t.me/aakhamidullin}{Telegram: @aakhamidullin} | 
    \href{https://github.com/Ahamidullin}{github.com/Ahamidullin} | 
    Москва
\end{center}


% ====== ОБРАЗОВАНИЕ ======
\section{ОБУЧЕНИЕ}
\textbf{НИУ ВШЭ, Москва} — Прикладная математика и информатика / Прикладной анализ данных, 3 курс \hfill 2023–2027

\section{ЦЕЛЬ}
Ищу стажировку/позицию в области Machine Learning, Deep Learning, Data Science, Recommender Systems. 

% ====== ОПЫТ ======
\section{ПРОФЕССИОНАЛЬНЫЙ ОПЫТ}
\textbf{S8 Capital} — Стажёр-аналитик \hfill Июнь–Август 2024
\begin{itemize}[left=0pt, itemsep=1pt]
    \item Анализировал данные 500K+ пользователей в Oracle: сегментация, таргетинг кампаний
    \item Проводил A/B-тесты (t-test, $\chi^2$-test), оценивал значимость через p-values и доверительные интервалы
    \item Применял bootstrap для проверки устойчивости результатов
    \item Автоматизировал расчёт метрик (конверсия, ARPU, retention) через Python/SQL в Jupyter Notebook
    \item Формировал аналитические отчёты с рекомендациями для продуктовой команды
\end{itemize}


% ====== ХАКАТОНЫ ======
\section{ХАКАТОНЫ}

\textbf{MTS True Tech Hack 2026} — ИИ-ассистент на базе OpenWebUI + MWS GPT API \hfill 2026
\begin{itemize}[left=0pt, itemsep=1pt]
    \item Deep Research агент: поиск (SearXNG) → скрейпинг (Jina) → LLM-синтез отчёта
    \item Голосовой ассистент: faster-whisper ASR + TTS, русский язык
    \item Система регистрации и онбординга пользователей; \href{https://github.com/hihihahovka/mts-hack}{GitHub}
\end{itemize}

\textbf{«Рекламатон»-2025} — Межвузовский хакатон по data-driven коммуникациям (ФКИ ВШЭ) \hfill 2025
\begin{itemize}[left=0pt, itemsep=1pt]
    \item Конструктор AI-персонажей: выбор архетипов, кастомизация стиля, диалог через LLM API
\end{itemize}

% \textbf{Хакатон по динамическому ценообразованию} — ФКН ВШЭ $\times$ VK Education \hfill 2026
% \begin{itemize}[left=0pt, itemsep=1pt]
%     \item Анализ ценовой эластичности, влияния промо-акций и складских остатков на спрос
% \end{itemize}

\textbf{HAHATHON} — Young\&\&Yandex $\times$ CSTATI \hfill 2026
\begin{itemize}[left=0pt, itemsep=1pt]
    \item «Яндекс Минус» — пародийный стриминг: антиплейлисты, случайные усложнения, юмористический UX (React)
    \item Голосовой ввод (Web Speech API), система регистрации пользователей
    \item \href{https://smehsmehsmeh.vercel.app}{smehsmehsmeh.vercel.app} | \href{https://github.com/hihihahovka/smehsmehsmeh}{GitHub} | \textbf{Победители}
\end{itemize}

% ====== ПРОЕКТЫ ======
\section{ПРОЕКТЫ}
\textbf{Fraud Detection System} | Python, ML, FastAPI, React | \href{https://github.com/Ahamidullin/Course-work}{GitHub}
\begin{itemize}[left=0pt, itemsep=1pt]
    \item Построил end-to-end ML pipeline для выявления мошеннических банковских транзакций на датасете из 1M записей с сильным дисбалансом классов (0.3\% fraud)
    \item Реализовал feature engineering: создание ratio-признаков (возраст/стаж проживания, скоринг/возраст), логарифмические трансформации velocity-метрик для сглаживания выбросов, бинаризация длительности сессий
    \item Разработал stacking ensemble из четырёх моделей (Logistic Regression, Random Forest, XGBoost как base learners + LightGBM meta-learner) с применением SMOTEENN для балансировки классов
    \item Оптимизировал гиперпараметры через RandomizedSearchCV и подобрал оптимальный threshold по F1-score на out-of-fold предсказаниях
    \item Задеплоил модель через FastAPI REST API с React-фронтендом для интерактивной real-time оценки риска транзакций
\end{itemize}

\vspace{2pt}
\textbf{Multimodal Meme Retrieval} | Python, CLIP, FAISS, Transformers | \href{https://github.com/Ahamidullin}{GitHub}
\begin{itemize}[left=0pt, itemsep=1pt]
    \item Разработал мультимодальную поисковую систему по базе 10K мемов: текстовый запрос (RU/EN) $\to$ релевантные изображения
    \item Построил пайплайн аннотаций: OCR (EasyOCR + PaddleOCR) $\to$ VQA-описания (Qwen2.5-VL) $\to$ NSFW-фильтрация
    \item Реализовал hybrid retrieval: multi-vector search (bge-m3), BM25, CLIP image search с RRF fusion
    \item Провёл эксперимент по выбору модели эмбеддингов (8 моделей), языковой маршрутизации и ablation study компонентов пайплайна
    \item \textbf{Результат:} Hit@5 с 32\% до 51\% (EN), с 2\% до 46\% (RU) на 100 размеченных запросах
\end{itemize}

\vspace{2pt}
\textbf{HSE.TNDR — Telegram-бот для знакомств}
\begin{itemize}[left=0pt, itemsep=1pt]
    \item Проектирование микросервисной архитектуры (бот, backend, notifications, API)
    \item Построение дашбордов метрик, customer development, тестирование рекомендательных алгоритмов
    \item \textbf{Результат:} рост DAU 300→1000+, достижение 15K MAU
\end{itemize}

\vspace{2pt}
\textbf{Другие проекты:}
\begin{itemize}[left=0pt, itemsep=1pt]
    \item EDA на Kaggle: построение гипотез, корреляционный анализ, визуализация
    \item Qt-приложение для визуализации космических объектов + unit-тесты (C++)
    \item C++ приложение для кластеризации небесных тел
\end{itemize}

% ====== НАВЫКИ ======
\section{НАВЫКИ}
\begin{itemize}[left=0pt, itemsep=1pt]
    \item \textbf{Языки программирования и инструменты:} Python, SQL, PostgreSQL, Git, Linux, Docker, C++
    \item \textbf{ML/DS:} Machine Learning, Deep Learning, Data Science, PyTorch, Scikit-learn, Pandas, Polars, Numpy, Scipy, Optuna, WandB
    \item \textbf{Визуализация и деплой:} Matplotlib, Seaborn, FastAPI
    \item \textbf{Языки:} Русский (родной), Английский (B2)
\end{itemize}

% ====== ДОПОЛНИТЕЛЬНО ======
\section{РЕЛЕВАНТНЫЕ КУРСЫ}
\begin{itemize}[left=0pt, itemsep=1pt]
    \item Machine Learning (Е. Соколов, ВШЭ) | \href{https://github.com/esokolov/ml-course-hse}{GitHub}
    \item Introduction to Deep Learning (ВШЭ) — нейронные сети, CNN, RNN, Transformers | \href{https://github.com/xiyori/intro-to-dl-hse}{GitHub}
    \item Practical Reinforcement Learning (ШАД, Яндекс) | \href{https://github.com/yandexdataschool/Practical_RL}{GitHub}
    \item Нейросетевые рекомендательные системы (ВШЭ) | \href{https://github.com/KhrylchenkoKirill/DeepRecSys}{GitHub}
\end{itemize}

\end{document}
